%% 自動生成（phys/gen_phys.py）。直接編集しないこと
\makeatletter
\phdefvol{2}{熱学}{エネルギーを使いこなす}{「なぜ冷える？」「なぜ結露する？」に、熱で答えよう}
\phdefface{1}{ふくらむ}
\phdefface{2}{めぐる}
\phdefface{3}{移る}
\phdefface{4}{ぶつかる}
\def\phFrontLine{エンジンも、エアコンも、発電所も。熱学の4つの顔で読みとく}
\phdefpart{1}{第1部}{ふくらむ：熱と第一法則}{（熱-1〜熱-4）}{1}
\phdefpart{2}{第2部}{めぐる：第二法則とサイクル}{（熱-5〜熱-10）}{2}
\phdefpart{3}{第3部}{移る：伝熱}{（熱-11〜熱-13）}{3}
\phdefpart{4}{第4部}{ぶつかる：ミクロから見た熱}{（熱-14〜熱-16）}{4}
\def\ph@whyA{冷たいコップのまわりに、なぜ水がつく？}\def\ph@whyrA{熱-9\hspace{0.8em}\bkpageno{152}}
\def\ph@whyB{高い山の上で、おかしの袋がふくらむのはなぜ？}\def\ph@whyrB{熱-3\hspace{0.8em}\bkpageno{48}}
\def\ph@whyC{お湯と水をまぜたあと、なぜ元の温度に分かれない？}\def\ph@whyrC{熱-6\hspace{0.8em}\bkpageno{100}}
\def\ph@whyD{二重窓にすると、なぜ部屋が冷えにくい？}\def\ph@whyrD{熱-11\hspace{0.8em}\bkpageno{188}}
\def\ph@whyE{ディーゼルエンジンは、なぜ火花なしで燃える？}\def\ph@whyrE{熱-7\hspace{0.8em}\bkpageno{118}}
\def\ph@whypage{312}
\def\phBackCatch{気になる「なぜ？」から、\\ページをめくってみよう}
\makeatother
